% !TEX TS-program = tectonic
% !TEX encoding = UTF-8 Unicode
\documentclass[11pt,a4paper]{article}

\usepackage{fontspec}
\usepackage{polyglossia}
\setdefaultlanguage{ukrainian}
\setotherlanguage{english}

\setmainfont{PT Serif}[Ligatures=TeX]
\setsansfont{PT Sans}[Ligatures=TeX]
\setmonofont{Menlo}[Scale=0.84]
\newfontfamily\cyrillicfont{PT Serif}[Ligatures=TeX]
\newfontfamily\cyrillicfontsf{PT Sans}[Ligatures=TeX]
\newfontfamily\cyrillicfonttt{Menlo}[Scale=0.84]

\usepackage[a4paper,top=2.25cm,bottom=2.30cm,left=2.35cm,right=2.35cm]{geometry}
\usepackage{amsmath}
\usepackage{amssymb}
\usepackage{booktabs}
\usepackage{tabularx}
\usepackage{longtable}
\usepackage{array}
\usepackage{enumitem}
\usepackage{xcolor}
\usepackage[most]{tcolorbox}
\usepackage{titlesec}
\usepackage{fancyhdr}
\usepackage[hidelinks]{hyperref}
\usepackage{microtype}
\usepackage{graphicx}
\usepackage{multicol}
\usepackage{listings}
\usepackage{tikz}
\usetikzlibrary{arrows.meta,positioning,shapes.geometric,fit}

\definecolor{accent}{HTML}{1F4E79}
\definecolor{accenttwo}{HTML}{2F6F5E}
\definecolor{soft}{HTML}{F2F5F8}
\definecolor{softgreen}{HTML}{F1F7F4}
\definecolor{softamber}{HTML}{FFF8E8}
\definecolor{softred}{HTML}{FFF3F1}
\definecolor{rule}{HTML}{C8D2DC}
\definecolor{warning}{HTML}{9A6700}
\definecolor{danger}{HTML}{A33A2B}

\hypersetup{
  colorlinks=true,
  linkcolor=accent,
  urlcolor=accent,
  citecolor=accent,
  pdftitle={When Does Schema Context Help? Повний план schema-aware статті},
  pdfauthor={Volodymyr Sabadosh and Vladyslav Kotsovsky}
}
\urlstyle{same}

\titleformat{\section}
  {\sffamily\Large\bfseries\color{accent}}{\thesection.}{0.6em}{}
\titleformat{\subsection}
  {\sffamily\large\bfseries}{\thesubsection.}{0.6em}{}
\titleformat{\subsubsection}
  {\sffamily\normalsize\bfseries}{\thesubsubsection.}{0.6em}{}
\titlespacing*{\section}{0pt}{18pt}{8pt}
\titlespacing*{\subsection}{0pt}{14pt}{5pt}
\titlespacing*{\subsubsection}{0pt}{11pt}{4pt}

\pagestyle{fancy}
\fancyhf{}
\renewcommand{\headrulewidth}{0.4pt}
\fancyhead[L]{\sffamily\footnotesize\color{accent}SchemaAwareArticlePlan}
\fancyhead[R]{\sffamily\footnotesize\color{accent}Catalog-grounded SQL auditing}
\fancyfoot[C]{\sffamily\footnotesize\thepage}

\newtcolorbox{keybox}[1][]{
  enhanced, breakable, arc=2pt, boxrule=0.8pt,
  colback=soft, colframe=accent!55,
  left=8pt, right=8pt, top=6pt, bottom=6pt,
  #1}

\newtcolorbox{thesisbox}[1][]{
  enhanced, breakable, arc=2pt, boxrule=0pt,
  colback=soft, frame hidden,
  borderline west={3pt}{0pt}{accent},
  left=10pt, right=8pt, top=6pt, bottom=6pt,
  #1}

\newtcolorbox{mustbox}[1][]{
  enhanced, breakable, arc=2pt, boxrule=0.8pt,
  colback=softgreen, colframe=accenttwo!65,
  title={\sffamily\bfseries Обов'язково для публікабельної версії},
  coltitle=black,
  left=8pt, right=8pt, top=5pt, bottom=6pt,
  #1}

\newtcolorbox{warningbox}[1][]{
  enhanced, breakable, arc=2pt, boxrule=0.8pt,
  colback=softamber, colframe=warning!65,
  title={\sffamily\bfseries Методологічне застереження},
  coltitle=black,
  left=8pt, right=8pt, top=5pt, bottom=6pt,
  #1}

\newtcolorbox{riskbox}[1][]{
  enhanced, breakable, arc=2pt, boxrule=0.8pt,
  colback=softred, colframe=danger!55,
  title={\sffamily\bfseries Ризик для claims},
  coltitle=black,
  left=8pt, right=8pt, top=5pt, bottom=6pt,
  #1}

\newcommand{\core}{\textcolor{accenttwo}{\textbf{CORE}}}
\newcommand{\secondary}{\textcolor{warning}{\textbf{SECONDARY}}}
\newcommand{\future}{\textcolor{danger}{\textbf{FUTURE}}}
\newcommand{\unknown}{\texttt{unknown}}
\newcommand{\presentstate}{\texttt{present}}
\newcommand{\absentstate}{\texttt{absent}}
\newcolumntype{Y}{>{\raggedright\arraybackslash}X}

\setlist[itemize]{leftmargin=1.4em,itemsep=2pt,topsep=4pt}
\setlist[enumerate]{leftmargin=1.6em,itemsep=3pt,topsep=4pt}
\renewcommand{\arraystretch}{1.22}
\setlength{\emergencystretch}{3em}

\lstdefinestyle{sqlstyle}{
  language=SQL,
  basicstyle=\ttfamily\small,
  keywordstyle=\color{accent}\bfseries,
  commentstyle=\color{accenttwo},
  stringstyle=\color{warning},
  showstringspaces=false,
  breaklines=true,
  frame=single,
  rulecolor=\color{rule},
  backgroundcolor=\color{soft},
  columns=fullflexible,
  xleftmargin=3pt,
  xrightmargin=3pt,
  aboveskip=5pt,
  belowskip=6pt
}

\lstdefinestyle{pythonstyle}{
  language=Python,
  basicstyle=\ttfamily\small,
  keywordstyle=\color{accent}\bfseries,
  commentstyle=\color{accenttwo},
  stringstyle=\color{warning},
  showstringspaces=false,
  breaklines=true,
  frame=single,
  rulecolor=\color{rule},
  backgroundcolor=\color{soft},
  columns=fullflexible,
  xleftmargin=3pt,
  xrightmargin=3pt,
  aboveskip=5pt,
  belowskip=6pt
}

\begin{document}

\begin{center}
{\sffamily\LARGE\bfseries\color{accent} Повний план окремої schema-aware статті}\\[5pt]
{\sffamily\large When Does Schema Context Help?\\
A Controlled Study of Catalog-Grounded SQL Auditing\\
for Text-to-SQL Benchmarks}\\[11pt]
{\footnotesize\sffamily
Робочий документ • Volodymyr Sabadosh, Vladyslav Kotsovsky • Uzhhorod National University\\
Пряме продовження AST-only роботи UJIT; окрема лінія від calibration/reconciliation article\\
Версія 1.0 • 17 серпня 2026}
\end{center}

\vspace{7pt}
\noindent\textcolor{rule}{\rule{\textwidth}{0.8pt}}
\vspace{9pt}

\begin{thesisbox}
\textbf{Головна теза.} AST показує структурну форму SQL, але не може самостійно визначити,
чи існує колонка, чи є вона nullable, чи задає GROUP BY функціональну залежність, чи
JOIN множить рядки та чи ORDER BY встановлює повний порядок. Нова робота має не просто
додати metadata до детектора, а \emph{контрольовано виміряти}, коли catalog grounding
пригнічує false positives, знаходить нові ризики або змушує систему чесно утриматися
від висновку.
\end{thesisbox}

\begin{keybox}
\textbf{Вердикт щодо публікабельності.} Сильна стаття можлива як дослідження
\emph{incremental value of schema context}: той самий SQL і той самий AST-rule family
порівнюються в режимах AST-only, declared schema, native catalog і consistency-checked
catalog. Варіант «ми додали PK/FK до лінтера» недостатній, оскільки SQLCheck уже
використовував контекст, а новий інструмент \texttt{sqlsure} уже перевіряє join cardinality
на Spider і BIRD.
\end{keybox}

\begin{mustbox}
\textbf{Публічні посилання, які мають бути в документі та майбутній статті:}
\begin{itemize}
\item Наша опублікована UJIT-робота:
  \href{https://doi.org/10.23939/ujit2026.01.135}
  {\textit{Automated Detection and Comparative Analysis of SQL Antipatterns in Text-to-SQL Datasets}}
  (\href{https://science.lpnu.ua/sites/default/files/journal-paper/2026/apr/42369/15sabadoshkocovskiy135-143.pdf}{прямий PDF}).
\item Новий споріднений інструмент, який уточнив дизайн цієї наступної роботи:
  \href{https://github.com/sqlsure/sqlsure}{\texttt{sqlsure}};
  аналізований frozen snapshot:
  \href{https://github.com/sqlsure/sqlsure/tree/d1b98e741f848ce7d756ee1d672e206d6ec39daa}{commit \texttt{d1b98e7}}.
\item Дослідницька система SQLCheck:
  \href{https://doi.org/10.1145/3318464.3389754}
  {\textit{SQLCheck: Automated Detection and Diagnosis of SQL Anti-Patterns}}.
\end{itemize}
\end{mustbox}

\tableofcontents
\clearpage

%==================================================================
\section{Executive summary і межі роботи}
%==================================================================

\subsection{Що саме будуємо}

До наявного dialect-aware AST analyzer додається спільний \texttt{SchemaContext},
який отримується двома шляхами:

\begin{enumerate}
\item із native database через catalog reflection;
\item із DDL/\texttt{sql\_context}: наявний materializer створює порожню БД, після чого
      використовується той самий reflection path.
\end{enumerate}

Rule engine залишається чистим: він не відкриває DB connection і не виконує SQL.
Він отримує AST, immutable context та повертає evidence-backed рішення. Запит не
запускається вдруге; execution analyzer продовжує виконувати його один раз у своєму
окремому pipeline step.

\subsection{Що є новим результатом}

\begin{itemize}
\item Контрольований C0--C3 ablation, а не повторний corpus census.
\item Schema facts із provenance і reliability state.
\item Scope/alias/CTE/derived-table lineage, потрібний для коректного зв'язування AST
      із колонками.
\item Тризначна семантика \presentstate/\absentstate/\unknown замість silent-safe.
\item Незалежний risk vector: correctness, performance, portability, maintainability.
\item Нова blind annotation set для disagreement cases.
\item Систематичне порівняння з \texttt{sqlsure} лише на спільних rule families.
\end{itemize}

\subsection{Що свідомо не входить}

\begin{itemize}
\item Повна NL--SQL semantic correctness evaluation: це належить окремій статті
      про calibration та reconciliation.
\item Автоматичне переписування SQL. Дозволена лише текстова remediation hint,
      підкріплена evidence.
\item Універсальний optimizer oracle для всіх діалектів.
\item Partition-aware optimization: Spider/BIRD SQLite не надають відповідного
      основного експериментального середовища.
\item Твердження, що відсутній FK доводить неправильність JOIN.
\end{itemize}

%==================================================================
\section{Чітка дельта до наших попередніх робіт}
%==================================================================

\subsection{UJIT 2026 як C0 baseline}

UJIT-стаття реалізувала 14 AST-based antipattern types, dialect-specific severity
profiles та corpus-level аналіз 122\,802 SQL-запитів у Spider, BIRD і Gretel.
Її ключове обмеження сформульовано прямо: exclusive reliance on AST-level static
analysis без column types, nullability, indexes, partitioning та execution-plan
evidence. У future work названо schema-aware detection, severity calibration і
multi-dimensional risk model.

\begin{table}[h]
\centering
\small
\begin{tabularx}{\textwidth}{>{\bfseries}p{3.2cm}YY}
\toprule
Вимір & UJIT & Нова schema-aware робота \\
\midrule
Основний об'єкт &
AST signatures і corpus profiles &
Зміна рішення одного rule family після catalog grounding \\
Контекст &
SQL AST + dialect profile &
AST + resolved lineage + schema facts + provenance \\
Вердикт &
Boolean flag + single severity &
\presentstate/\absentstate/\unknown + risk vector \\
Оцінювання &
Prevalence і порівняння із C++ SqlCheck &
Matched ablation, blind disagreement review, \texttt{sqlsure} baseline \\
Новизна &
Corpus-scale AST antipattern audit &
Measured incremental value and limits of schema context \\
\bottomrule
\end{tabularx}
\caption{Дельта до \href{https://doi.org/10.23939/ujit2026.01.135}{UJIT 2026}.}
\end{table}

\subsection{e-Informatica як infrastructure baseline}

e-Informatica manuscript уже має:

\begin{itemize}
\item schema validation analyzer;
\item \texttt{DbManager} і adapter protocol;
\item materialization DB із DDL, якщо native DB відсутня;
\item column type normalization, PK/FK reflection, referential-integrity checks;
\item syntax, execution, AST-antipattern та LLM-semantic analyzers.
\end{itemize}

Але schema analyzer та query antipattern analyzer працюють паралельно. Нова робота
з'єднує їх: schema facts стають входом для query-level rules. Старі schema totals,
execution totals, 17 Cartesian cases і 263 semantic labels не подаються як нові
результати.

\subsection{План calibration та reconciliation залишається окремим}

Calibration article оцінює semantic defect prevalence та recall LLM committee, а також
item-level reconciliation із SQLDriller. Поточна стаття оцінює structural
schema-conditioned propositions. Вона не використовує старі 263 labels як ground
truth і не намагається дати новий загальний «відсоток неправильних gold SQL».

\begin{warningbox}
У майбутній статті потрібен окремий підрозділ \emph{Relation to our prior work}.
Без нього рецензент може сприйняти C0 і повторне використання корпусів як salami
slicing. Захистом є нова одиниця аналізу: \emph{paired rule decision under changing
context}, а не повторна таблиця prevalence.
\end{warningbox}

%==================================================================
\section{SQLCheck і sqlsure: що вже зроблено іншими}
%==================================================================

\subsection{SQLCheck}

SQLCheck SIGMOD 2020 сформулював pipeline context formation, anti-pattern detection,
ranking і repair для database applications. Тому не можна заявляти «first
schema-aware SQL antipattern detector». Наша дельта: gold-query benchmark auditing,
ідентичний AST-only/schema-aware ablation, provenance, abstention і незалежна
емпірична оцінка.

Публічно відтворюваний C++ \href{https://github.com/jarulraj/sqlcheck}{SqlCheck}
є іншим artifact і вже використовувався в UJIT як baseline. Його можна залишити
secondary baseline лише для overlapping syntax rules; він не є центральним
schema-aware конкурентом.

\subsection{Що реально реалізує sqlsure}

Станом на 17 серпня 2026 року \texttt{sqlsure} є open-source artifact v0.1.1, а не
рецензованою статтею. Його core:

\begin{itemize}
\item Python + SQLGlot;
\item semantic model із table grain, measures, sensitive columns і declared joins;
\item SQLite \texttt{PRAGMA} та PostgreSQL/MySQL \texttt{information\_schema}
      introspection для PK/FK;
\item дев'ять rules: FANOUT, CHASM, ADDITIVITY, SEMI\_ADDITIVE, WEIGHTED\_AVG,
      JOIN\_KEY, UNDECLARED\_JOIN, CROSS\_JOIN, SENSITIVE\_COLUMN.
\end{itemize}

Catalog автоматично дає grain і join cardinality. Additivity, semi-additive dimension
і sensitive policy потребують semantic metadata або ручної декларації. Інструмент
не читає rows, не перевіряє RI, не є data-aware й не знає NL intent.

\subsection{Важливі обмеження sqlsure для fair comparison}

Аналіз frozen snapshot \texttt{d1b98e7} показує:

\begin{itemize}
\item parser scope-aware, але не повністю lineage-aware;
\item CTE output lineage не переноситься в outer scopes;
\item derived-table joins можуть бути пропущені;
\item partial composite-key match може помилково пройти як valid join;
\item \texttt{ON TRUE} та \texttt{ON 1=1} мають predicate і не позначаються CROSS\_JOIN;
\item порожній список violations іноді означає «не вдалося атрибутувати», а не
      доведене safe;
\item Spider і BIRD audit scripts викликають extraction layer та окремо класифікують
      join pairs із назв колонок; це не тотожне запуску всіх дев'яти \texttt{check()} rules;
\item «zero false alarms» базується на retrospective manual review candidates і
      не є blind evaluation всього rule set.
\end{itemize}

\begin{keybox}
\textbf{Науково коректна позиція.} \texttt{sqlsure} є корисним, пізнішим і вузьким
semantic-model checker. Він сильніший за UJIT у cardinality/additivity domain, але
не перекриває UJIT taxonomy або e-Informatica pipeline. Нова стаття повинна
цитувати artifact, pin exact commit і використовувати його як baseline, а не
знецінювати.
\end{keybox}

\subsection{Хронологія і priority}

\begin{itemize}
\item UJIT article received: 31 березня 2026.
\item UJIT article accepted: 22 квітня 2026.
\item \texttt{sqlsure} repository created: 2 липня 2026.
\end{itemize}

Отже, UJIT має чіткий хронологічний priority для опублікованого corpus-scale
SQLGlot AST audit. \texttt{sqlsure} впливає на позиціонування наступної
schema-aware статті, але не ретроспективно на новизну UJIT.

%==================================================================
\section{Research gap, RQs і гіпотези}
%==================================================================

\subsection{Точний gap}

\begin{thesisbox}
У знайденій літературі та artifacts немає peer-reviewed контрольованого дослідження,
яке на одному й тому самому rule engine із незалежними labels вимірює incremental
value of declared, native та consistency-checked schema facts для аудиту reference
SQL, включно з explicit abstention за неповного metadata.
\end{thesisbox}

\subsection{Research questions}

\begin{description}[leftmargin=1.5cm,labelwidth=1.0cm,style=nextline]
\item[RQ1] Наскільки повними, узгодженими й надійними є names, types, nullability,
PK/FK, uniqueness та indexes у Spider, BIRD і Gretel?
\item[RQ2] Які AST-only findings context підтверджує, пригнічує, додає або переводить
у \unknown?
\item[RQ3] Як context змінює precision, false-positive rate, review yield і
risk-vector assignments?
\item[RQ4] Які rule families найбільше виграють від nullability, key, cardinality,
type, index і lineage facts?
\item[RQ5] Які підтверджені cases унікальні для нашого analyzer, \texttt{sqlsure},
DBMS resolution та execution check?
\item[RQ6] Як acquisition path --- native database проти DDL-materialized schema ---
впливає на coverage, abstention і throughput?
\end{description}

\subsection{Pre-specified hypotheses}

\begin{itemize}
\item \textbf{H1:} context-dependent rules матимуть вищу precision у C1--C3, ніж у C0.
\item \textbf{H2:} syntax-only negative-control rules не зміняться між conditions.
\item \textbf{H3:} найбільше false-positive suppression дасть nullable NOT IN.
\item \textbf{H4:} найбільше нових correctness-risk findings дадуть join-key mismatch,
      fan-out/chasm і potential multi-row scalar subquery.
\item \textbf{H5:} Gretel матиме більше \unknown через неповний або неоднорідний DDL,
      але materialization дозволить застосувати declaration-based rules масштабно.
\item \textbf{H6:} published metadata, native catalog і consistency checks на даних
      дадуть не тотожні rule decisions.
\item \textbf{H7:} lineage-aware аналіз матиме вищу applicable coverage на CTE та
      derived-table queries, ніж \texttt{sqlsure}.
\end{itemize}

%==================================================================
\section{Архітектура запропонованої системи}
%==================================================================

\subsection{Data flow}

\begin{center}
\begin{tikzpicture}[
  node distance=7mm and 9mm,
  every node/.style={font=\sffamily\small},
  box/.style={draw=rule, rounded corners=2pt, align=center, minimum height=9mm,
              text width=30mm, fill=soft},
  corebox/.style={draw=accent!60, rounded corners=2pt, align=center,
                  minimum height=10mm, text width=34mm, fill=softgreen},
  arrow/.style={-{Latex[length=2mm]}, thick, draw=accent!70}
]
\node[box] (sql) {Gold SQL\\+ dialect};
\node[box, below=of sql] (ast) {SQLGlot AST};
\node[box, right=of sql] (native) {Native DB};
\node[box, right=of ast] (ddl) {DDL / sql\_context};
\node[box, right=of ddl] (materialize) {Existing DB\\materializer};
\node[box, below=of native] (catalog) {Catalog reflection};
\node[corebox, below=of catalog] (context) {Immutable\\SchemaContext};
\node[corebox, below=of ast] (resolver) {Scope, alias\\and lineage resolver};
\node[corebox, below=of resolver] (rules) {Tri-state\\rule engine};
\node[box, right=of rules] (evidence) {Evidence + provenance\\+ risk vector};
\node[box, below=of evidence] (outputs) {JSONL / DuckDB\\reports};

\draw[arrow] (sql) -- (ast);
\draw[arrow] (native) -- (catalog);
\draw[arrow] (ddl) -- (materialize);
\draw[arrow] (materialize) |- (catalog);
\draw[arrow] (catalog) -- (context);
\draw[arrow] (ast) -- (resolver);
\draw[arrow] (context) -| (resolver);
\draw[arrow] (resolver) -- (rules);
\draw[arrow] (rules) -- (evidence);
\draw[arrow] (evidence) -- (outputs);
\end{tikzpicture}
\end{center}

\subsection{Не дублювати materialization}

Наявний lifecycle уже створює database/schema з DDL за відсутності native DB.
Нова реалізація не переносить цей код у detector. На analyzer boundary:

\begin{enumerate}
\item normalizer отримує schema/DDL;
\item adapter canonicalizes DDL і формує stable schema identity;
\item DB створюється лише за потреби;
\item \texttt{DbManager} повертає tables і table info;
\item \texttt{SchemaContextBuilder} freezes facts і кешує context за \texttt{db\_id}
      або canonical schema hash;
\item pure detector використовує context без I/O.
\end{enumerate}

Для Gretel однакові DDL мають дедуплікуватися за canonical hash. Порожня materialized
DB дозволяє names/types/nullability/PK/FK/unique/index rules, але не data-derived
claims.

\subsection{Модель контексту}

\begin{lstlisting}[style=pythonstyle]
@dataclass(frozen=True)
class ColumnFact:
    name: str
    type_raw: str
    type_family: str
    nullable: bool | None
    primary_key: bool
    unique: bool | None
    provenance: FactProvenance

@dataclass(frozen=True)
class TableFact:
    name: str
    columns: Mapping[str, ColumnFact]
    primary_key: tuple[str, ...]
    unique_keys: tuple[tuple[str, ...], ...]
    foreign_keys: tuple[ForeignKeyFact, ...]
    indexes: tuple[IndexFact, ...]

@dataclass(frozen=True)
class SchemaContext:
    db_id: str
    dialect: str
    source: str
    tables: Mapping[str, TableFact]
\end{lstlisting}

\subsection{Provenance і reliability}

\begin{table}[h]
\centering
\small
\begin{tabularx}{\textwidth}{>{\ttfamily}p{4.1cm}YY}
\toprule
\normalfont\textbf{Provenance} & \textbf{Що означає} & \textbf{Дозволені висновки} \\
\midrule
benchmark\_metadata &
Факт узято з tables.json або іншого benchmark manifest &
Declaration-based; missing fact не є доказом відсутності \\
\shortstack[l]{materialized\_from\\ddl} &
Catalog створено з DDL без population data &
Names, types, declared keys, nullability, indexes \\
native\_catalog &
Факт прочитано з shipped DB catalog &
Catalog-level reasoning; ще не доказ data consistency \\
\shortstack[l]{data\_consistency\\checked} &
Declared constraint не суперечить поточному instance &
Підвищена довіра лише для цього snapshot; не універсальний доказ \\
observed\_only &
Uniqueness/null absence спостережено в rows, але не declared &
Supporting evidence; не підмінює logical constraint \\
\bottomrule
\end{tabularx}
\caption{Джерела schema facts та межі інтерпретації.}
\end{table}

\subsection{Lineage resolver}

Resolver повинен:

\begin{itemize}
\item будувати lexical scope для кожного SELECT;
\item відображати alias у base table;
\item вирішувати qualified та unqualified columns;
\item виявляти ambiguity, а не вгадувати owner;
\item переносити output-column lineage через CTE та derived tables;
\item підтримувати projection aliases, \texttt{USING}, composite keys та nested scopes;
\item позначати unsupported expression lineage як \unknown.
\end{itemize}

Це важливіше за додавання десятків rules: помилкова атрибуція колонки робить будь-яке
schema-aware правило недостовірним.

\subsection{RuleDecision і risk vector}

\begin{lstlisting}[style=pythonstyle]
@dataclass(frozen=True)
class RuleDecision:
    rule_id: str
    state: Literal["present", "absent", "unknown"]
    evidence: tuple[Evidence, ...]
    provenance: tuple[FactProvenance, ...]
    correctness: RiskLevel | None
    performance: RiskLevel | None
    portability: RiskLevel | None
    maintainability: RiskLevel | None
    remediation_hint: str | None
\end{lstlisting}

Legacy Critical/High/Medium/Low залишається лише для backward compatibility. У
науковому аналізі не слід зводити correctness і performance до одного числа без
емпірично валідованих ваг.

%==================================================================
\section{Які наявні rules виграють від schema context}
%==================================================================

\begin{longtable}{>{\raggedright\arraybackslash}p{3.5cm}
                  >{\raggedright\arraybackslash}p{2.7cm}
                  >{\raggedright\arraybackslash}p{7.6cm}}
\toprule
\textbf{Наявний rule} & \textbf{Виграш} & \textbf{Schema-aware refinement} \\
\midrule
\endfirsthead
\toprule
\textbf{Наявний rule} & \textbf{Виграш} & \textbf{Schema-aware refinement} \\
\midrule
\endhead
NOT IN nullable & Дуже високий &
Визначити nullability projected expression, включно з outer-join/CASE lineage;
не flag усі NOT IN автоматично. \\
Missing GROUP BY & Високий &
Перевірити PK, composite/UNIQUE NOT NULL keys та functional dependencies;
розділити correctness і portability. \\
LIMIT/OFFSET ordering & Високий &
Не лише наявність ORDER BY, а total order відносно result grain. \\
Function in WHERE & Високий за наявності metadata &
Типи, implicit conversion, ordinary/expression indexes; EXPLAIN subset validation. \\
Cartesian product & Середній &
Краща alias ownership, FK path evidence, wrong-table/self-reference diagnosis;
сам Cartesian факт переважно AST-based. \\
Redundant DISTINCT & Середній--високий &
PK/unique/FD дозволяють довести, що output уже унікальний. \\
SELECT * & Середній для severity &
Розгорнути projection; BLOB/JSON/wide columns підвищують resource risk;
sensitive metadata може додати policy risk. \\
Leading wildcard LIKE & Середній для performance &
Index/collation/type metadata; без індексу не стверджувати «index disabled». \\
Correlated subquery & Середній &
Lineage підтверджує реальну correlation; index metadata уточнює performance risk. \\
NULL comparison & Мінімальний &
\texttt{= NULL} є некоректною predicate form незалежно від schema. \\
Unsafe UPDATE/DELETE & Мінімальний &
Відсутність WHERE визначається AST; schema не змінює основний verdict. \\
Columns in EXISTS & Мінімальний &
Projection усередині EXISTS семантично ігнорується; AST достатньо. \\
\bottomrule
\caption{Вплив schema context на опублікований UJIT rule set.}
\end{longtable}

\begin{mustbox}
Primary article scope: поглибити шість наявних families --- nullable NOT IN,
Missing GROUP BY, LIMIT/OFFSET total order, Function in WHERE, schema-grounded
Cartesian/join diagnosis і key-aware Redundant DISTINCT. SELECT *, leading wildcard
і correlated subquery залишити secondary analyses.
\end{mustbox}

%==================================================================
\section{Primary schema-aware rules: формальні ідеї та приклади}
%==================================================================

\subsection{R1: nullable NOT IN}

\textbf{Required facts:} lineage projected subquery expression, declared nullability,
null-introducing operators.

\begin{itemize}
\item \presentstate: expression може бути NULL за declared schema або expression
      semantics.
\item \absentstate: expression доведено NOT NULL у цьому scope.
\item \unknown: unresolved lineage або missing nullability fact.
\end{itemize}

\begin{lstlisting}[style=sqlstyle]
SELECT u.id
FROM users AS u
WHERE u.id NOT IN (
  SELECT b.user_id
  FROM blocked_users AS b
);
\end{lstlisting}

Якщо \texttt{blocked\_users.user\_id} nullable, один NULL може зробити predicate
UNKNOWN для всіх candidate rows. Якщо колонка \texttt{NOT NULL}, AST-only flag
пригнічується. Якщо projection є CASE або походить із nullable side outer join,
простого column-level NOT NULL недостатньо.

\subsection{R2: functional-dependency-aware Missing GROUP BY}

\begin{lstlisting}[style=sqlstyle]
CREATE TABLE customers (
  id INTEGER PRIMARY KEY,
  name TEXT NOT NULL
);

SELECT c.id, c.name, COUNT(o.id)
FROM customers AS c
LEFT JOIN orders AS o ON o.customer_id = c.id
GROUP BY c.id;
\end{lstlisting}

AST-only detector бачить \texttt{c.name} поза GROUP BY. Schema-aware detector
доводить:
\[
\texttt{customers.id} \rightarrow \texttt{customers.name}.
\]
Correctness risk стає absent; portability може залишитися dialect-dependent.

Контрприклад:

\begin{lstlisting}[style=sqlstyle]
SELECT country, city, COUNT(*)
FROM customers
GROUP BY country;
\end{lstlisting}

Якщо \texttt{country} не покриває PK або UNIQUE NOT NULL key, \texttt{city} не є
функціонально залежною і SQLite може вибрати довільне значення. Для nullable UNIQUE
потрібна обережність: кілька NULL можуть утворити одну групу. Для composite key
потрібні всі key columns.

\subsection{R3: total-order-aware LIMIT/OFFSET}

\begin{lstlisting}[style=sqlstyle]
SELECT id, score
FROM candidates
ORDER BY score DESC
LIMIT 10;
\end{lstlisting}

Наявність ORDER BY ще не гарантує deterministic top-10: ties за \texttt{score}
дозволяють різний subset. Якщо \texttt{score} не unique, recommendation:
\texttt{ORDER BY score DESC, id}. Rule має враховувати result grain після GROUP BY,
DISTINCT і joins; якщо grain не виведено, verdict \unknown.

\subsection{R4: type/index-aware SARGability}

\begin{lstlisting}[style=sqlstyle]
SELECT *
FROM events
WHERE CAST(user_id AS TEXT) = '42';
\end{lstlisting}

AST-only rule бачить function/cast. Schema-aware version перевіряє:

\begin{itemize}
\item actual type family \texttt{user\_id};
\item implicit/explicit conversion direction;
\item ordinary або expression index;
\item dialect-specific coercion;
\item plan evidence на підтримуваній subset через \texttt{EXPLAIN QUERY PLAN}.
\end{itemize}

Без індексу не слід стверджувати, що function «заблокувала index». Тоді може
залишитися computation/portability risk, але не підтверджений index-access defect.

\subsection{R5: schema-grounded Cartesian/join diagnosis}

AST уже знаходить comma joins, \texttt{ON TRUE}, one-sided predicates та missing
ON. Schema додає:

\begin{itemize}
\item точну ownership колонок;
\item expected FK path як evidence;
\item distinction між wrong key, absent predicate та unknown undeclared relation;
\item diagnosis wrong-table references.
\end{itemize}

Відсутній FK не перетворюється на error: це \unknown або unmodeled warning.

\subsection{R6: key-aware Redundant DISTINCT}

DISTINCT можна довести redundant, якщо projected columns містять result key і joins
не розмножують цей grain. Навпаки, наявність PK базової таблиці не достатня після
one-to-many join. Потрібні lineage і cardinality propagation.

%==================================================================
\section{Три нові rules, потрібні для Text-to-SQL}
%==================================================================

\subsection{N1: aggregate after fan-out або chasm}

\begin{lstlisting}[style=sqlstyle]
SELECT c.id, SUM(c.credit_limit)
FROM customers AS c
JOIN orders AS o ON o.customer_id = c.id
GROUP BY c.id;
\end{lstlisting}

Якщо \texttt{customers.id -> orders} є one-to-many, base-row
\texttt{credit\_limit} повторюється для кожного order. Schema facts доводять row
multiplication; intent щодо measure усе ще може вимагати review.

\begin{itemize}
\item \textbf{FANOUT:} aggregate над measure/column із grain, який повторюється join.
\item \textbf{CHASM:} дві або більше fan-out branches множать одна одну.
\item \textbf{Не робити:} вважати \texttt{SUM(DISTINCT amount)} універсально безпечним,
      бо різні rows можуть мати однаковий amount.
\end{itemize}

Цей rule напряму порівнюється із \texttt{sqlsure} FANOUT/CHASM, але наш експеримент
додає lineage, provenance, abstention і independent labels.

\subsection{N2: join-key mismatch}

\begin{lstlisting}[style=sqlstyle]
SELECT *
FROM orders AS o
JOIN customers AS c
  ON o.sales_rep_id = c.customer_id;
\end{lstlisting}

Якщо schema декларує \texttt{orders.customer\_id -> customers.customer\_id},
використаний key суперечить declared relationship. Critical correctness risk
допустимий лише за позитивного evidence. Якщо relationship не declared, система
повертає \unknown, а не автоматичний error.

Composite FK вважається matched тільки тоді, коли покрито весь key; один правильний
pair не повинен маскувати другий неправильний predicate.

\subsection{N3: potential multi-row scalar subquery}

\begin{lstlisting}[style=sqlstyle]
SELECT p.name,
       (SELECT r.rating
        FROM reviews AS r
        WHERE r.product_id = p.id) AS rating
FROM products AS p;
\end{lstlisting}

Якщо \texttt{reviews.product\_id} не unique і немає aggregate/LIMIT із
детермінованим order, scalar subquery може повернути кілька rows. PostgreSQL
завершиться error; SQLite може використати один рядок, створюючи portability та
determinism risk.

\subsection{Optional: outer-join null rejection}

\begin{lstlisting}[style=sqlstyle]
SELECT c.id
FROM customers AS c
LEFT JOIN orders AS o ON o.customer_id = c.id
WHERE o.status = 'paid';
\end{lstlisting}

WHERE predicate відкидає NULL-extended rows і фактично робить LEFT JOIN подібним до
INNER JOIN. Це важливий Text-to-SQL risk, але intent-dependent; додавати в primary
study лише після pilot.

%==================================================================
\section{Multi-dimensional risk model}
%==================================================================

\subsection{Чому single severity недостатня}

UJIT справедливо використовував Critical/High/Medium/Low як operational
prioritization. Schema context дозволяє розділити незалежні наслідки:

\begin{table}[h]
\centering
\small
\begin{tabularx}{\textwidth}{>{\bfseries}p{3cm}YY}
\toprule
Вісь & Питання & Приклади \\
\midrule
Correctness &
Чи може query повернути логічно неправильні rows/value? &
nullable NOT IN, fan-out SUM, wrong join key \\
Performance &
Чи є evidence про scan, row explosion або expensive projection? &
indexed function predicate, wide SELECT *, correlated execution \\
Portability &
Чи залежить поведінка від dialect/coercion? &
bare GROUP BY column, multi-row scalar subquery, type coercion \\
Maintainability &
Чи підвищує форма query ризик майбутньої помилки? &
SELECT *, undeclared join path, ambiguous ownership \\
\bottomrule
\end{tabularx}
\caption{Незалежні risk dimensions.}
\end{table}

Для кожної осі зберігати evidence-backed level або \unknown. Не публікувати
композитний quality score, доки ваги не валідовані людьми або downstream outcome.

%==================================================================
\section{Експериментальний протокол}
%==================================================================

\subsection{Corpora і версії}

\begin{itemize}
\item \textbf{Spider 1.0}: primary database-backed corpus, усі доступні partitions;
      exact source checksum у manifest.
\item \textbf{BIRD}: відтворити \texttt{dev\_20240627}, який використовував
      \texttt{sqlsure}, і окремо провести main evaluation на актуальному corrected
      snapshot, pinned до дати запуску.
\item \textbf{Gretel synthetic Text-to-SQL}: DDL-derived external-validity та
      scalability corpus; snapshot/commit, DDL hash і deduplication обов'язкові.
\end{itemize}

Spider і BIRD є primary study. Gretel звітується окремо; його synthetic SQL не
називається human gold SQL, а результати не pool-яться з database-backed corpora.

\subsection{Conditions}

\begin{description}[leftmargin=1.5cm,labelwidth=1.1cm]
\item[C0] AST-only UJIT-compatible mode.
\item[C1] AST + declared schema facts: benchmark metadata або DDL; за потреби
          acquisition через materialized empty DB.
\item[C2] AST + native live catalog shipped database.
\item[C3] AST + catalog constraints, consistency-checked against populated instance.
\end{description}

C3 не «доводить» general constraint: він лише показує, що current snapshot не
суперечить declaration. Для Gretel primary comparison C0/C1; reflection із
materialized DB є acquisition path для C1, а не незалежним source condition.

\subsection{Development/holdout discipline}

\begin{enumerate}
\item Розробляти rule logic на synthetic matched fixtures і Spider Train.
\item Freeze rule specification, severities/risk mapping і code commit.
\item Evaluate на Spider Dev/Test, corrected BIRD Dev і Gretel snapshot.
\item Known corrected BIRD/sqlsure cases використовувати лише як external validation
      після freeze або чесно позначити retrospective subset.
\end{enumerate}

Це зменшує corpus-guided tuning, яке послаблює як UJIT regression construction, так і
\texttt{sqlsure} retrospective claims.

\subsection{Baselines}

\begin{enumerate}
\item \textbf{C0 current analyzer}: exact published rule behavior.
\item \textbf{sqlsure v0.1.1, commit d1b98e7}: запуск public \texttt{check()} API.
\item \textbf{sqlsure audit scripts}: окремий результат, не змішувати з checker.
\item \textbf{DBMS prepare/resolution}: identifier/type hard errors.
\item \textbf{C++ SqlCheck v1.3}: secondary baseline для відтворюваних overlapping
      syntax rules.
\end{enumerate}

\begin{warningbox}
Не рахувати один глобальний F1 між systems із різними taxonomies. Порівняння
\texttt{sqlsure} обмежити CROSS\_JOIN, JOIN\_KEY, FANOUT, CHASM та clearly mapped
propositions. Additivity і sensitive policy не можна оцінювати без однакового
semantic metadata.
\end{warningbox}

\subsection{Matched synthetic tests як causal sanity check}

SQL залишається ідентичним, змінюється один schema fact:

\begin{itemize}
\item nullable $\leftrightarrow$ NOT NULL;
\item unique $\leftrightarrow$ non-unique;
\item FK present $\leftrightarrow$ absent;
\item one-to-one $\leftrightarrow$ one-to-many;
\item type-compatible $\leftrightarrow$ incompatible;
\item index absent $\leftrightarrow$ ordinary/expression index.
\end{itemize}

Syntax-only control rules повинні залишатися інваріантними. Якщо вони змінюються,
context integration має leakage або implementation bug.

\subsection{Blind human evaluation}

Нова annotation set формується зі страт:

\begin{itemize}
\item C0 present $\rightarrow$ C1/C2 absent;
\item C0 absent $\rightarrow$ C1/C2 present;
\item C1 $\neq$ C2 або C2 $\neq$ C3;
\item proposed-only та sqlsure-only disagreements;
\item random unflagged controls;
\item representative \unknown cases.
\end{itemize}

Два незалежні SQL/DB reviewers, blind до detector і condition, потім adjudication.
Label schema:

\begin{enumerate}
\item підтверджена query-level structural проблема;
\item valid query, але evidence-backed portability/performance risk;
\item schema/metadata defect;
\item intentional/acceptable construction;
\item unverifiable with available evidence.
\end{enumerate}

Початковий pilot визначає prevalence positives у strata; фінальний sample size
обирається через power/precision calculation. Реалістичний орієнтир для focused
rule-family claims --- 400--600 reviewed cases, але рідкісні classes можуть
потребувати census усіх findings.

\subsection{Metrics}

\begin{itemize}
\item context acquisition coverage і parse/resolve coverage;
\item rule applicability та abstention rate;
\item retained, suppressed, added і unknown migrations;
\item finding-level і query-level precision/recall/F1;
\item false-positive suppression і incremental true positives;
\item risk-vector migration;
\item overlap/Jaccard із baselines на matched rules;
\item inter-annotator agreement ($\kappa$ або Krippendorff's $\alpha$);
\item review time, throughput, schema-load latency, memory;
\item per-database clustering і concentration of findings.
\end{itemize}

Paired binary outcomes: McNemar. Confidence intervals: clustered bootstrap за
\texttt{db\_id}. Для stratified manual sample: sampling weights. Не трактувати
queries з однієї DB як незалежні observations.

\subsection{Performance subset}

Для Function-in-WHERE, leading wildcard і SELECT-*:

\begin{enumerate}
\item зібрати reliable index/type facts;
\item запустити \texttt{EXPLAIN QUERY PLAN} на підтримуваних SQLite cases;
\item відділити logical finding від plan-confirmed evidence;
\item не переносити SQLite plan conclusion на PostgreSQL без окремого experiment.
\end{enumerate}

%==================================================================
\section{Schema reliability і query-impact graph}
%==================================================================

\subsection{Навіщо перевіряти джерело фактів}

Published benchmark metadata, materialized DDL і native catalog можуть розходитись.
Пропонований audit повинен звітувати:

\begin{itemize}
\item missing/extra tables і columns;
\item nullability disagreements;
\item PK/FK/unique/index disagreements;
\item declared constraints, порушені current data;
\item unsupported або failed DDL materialization;
\item частку query references, які вдалося повністю resolve.
\end{itemize}

\subsection{Schema-to-query impact graph}

Нова корисна одиниця:
\[
\text{schema fact or defect} \longrightarrow
\{\text{dependent gold queries and rule decisions}\}.
\]

Наприклад, один missing FK може породити десятки unbacked-join candidates.
Замість «30 помилок SQL» треба показати «один schema defect, що впливає на 30
queries». Це запобігає double counting і напряму з'єднує e-Informatica schema
analyzer із query-level article.

%==================================================================
\section{Evidence records, counterexamples і відмова від auto-repair}
%==================================================================

\subsection{Evidence bundle}

Кожний finding повинен містити:

\begin{itemize}
\item rule id і precise AST location;
\item resolved table/column lineage;
\item schema facts і provenance;
\item reason for \presentstate/\absentstate/\unknown;
\item risk vector;
\item short human-readable remediation hint;
\item analyzer/rule/schema versions.
\end{itemize}

\subsection{Мінімальні witnesses}

Для кількох correctness rules можна створити невеликі counterexample instances:

\begin{itemize}
\item один NULL для NOT IN;
\item два child rows для fan-out;
\item ties для non-total ORDER BY;
\item два subquery rows для scalar-subquery risk.
\end{itemize}

Witness доводить, що risk може матеріалізуватися, але не доводить, що gold SQL не
відповідає конкретному NL question.

\subsection{Чому не робимо Repair}

Автоматичні rewrites для NOT IN, fan-out, GROUP BY або ORDER BY можуть змінити intent.
У primary scope:

\begin{itemize}
\item detector не змінює SQL;
\item повертає remediation hint;
\item safe rewrite на кшталт \texttt{= NULL -> IS NULL} залишається future extension;
\item semantic equivalence checking і repair evaluation потребують окремого study.
\end{itemize}

%==================================================================
\section{Implementation blueprint}
%==================================================================

\subsection{Existing integration points}

\begin{itemize}
\item \texttt{src/text2sql\_pipeline/analyzers/query\_antipattern/antipattern\_detector.py}
\item \texttt{.../query\_antipattern/query\_antipattern\_analyzer.py}
\item \texttt{.../query\_antipattern/antipattern\_registry.py}
\item \texttt{.../query\_antipattern/metrics.py}
\item \texttt{src/text2sql\_pipeline/db/manager.py}
\item \texttt{src/text2sql\_pipeline/db/adapters/base/protocol.py}
\item \texttt{.../schema\_validation/schema\_validation\_analyzer.py}
\item \texttt{src/text2sql\_pipeline/output/sinks/duckdb.py}
\end{itemize}

\subsection{Proposed modules}

\begin{itemize}
\item \texttt{schema\_context.py}: frozen facts, provenance, reliability.
\item \texttt{schema\_context\_builder.py}: native/DDL materialized acquisition і cache.
\item \texttt{scope\_resolver.py}: alias, CTE, derived lineage.
\item \texttt{rule\_decision.py}: tri-state evidence та risk vector.
\item \texttt{schema\_aware\_rules.py}: primary schema-dependent families.
\item \texttt{explain\_evidence.py}: optional supported-plan validation.
\end{itemize}

\subsection{Backward-compatible API}

\begin{lstlisting}[style=pythonstyle]
def detect_antipatterns(
    sql: str,
    dialect: str = "sqlite",
    config=None,
    penalties=None,
    schema_context: SchemaContext | None = None,
) -> QueryAntipatternFeatures:
    ...
\end{lstlisting}

Якщо context відсутній, C0 behavior зберігається. Schema-required rules не повинні
мовчки робити AST guess: залежно від pre-specified policy вони повертають \unknown
або не-applicable decision, який окремо рахується.

\subsection{Tests}

\begin{itemize}
\item усі наявні AST-only regression tests;
\item matched schema fixtures;
\item native DB vs DDL-materialized acquisition equivalence;
\item cache by \texttt{db\_id}/schema hash;
\item aliases, ambiguity, CTE shadowing, derived tables;
\item composite PK/FK/unique keys;
\item outer joins і null introduction;
\item SQLite affinity та quoted identifiers;
\item evidence serialization і DuckDB migrations;
\item no-row database не є evidence про data null/uniqueness.
\end{itemize}

\subsection{Implementation phases}

\begin{enumerate}
\item \textbf{P0 --- Specification freeze:} rule predicates, states, risk dimensions,
      datasets, baselines, pilot protocol.
\item \textbf{P1 --- Context foundation:} immutable facts, provenance, index reflection,
      shared cache, both acquisition paths.
\item \textbf{P2 --- Resolver:} scopes, aliases, CTE/derived lineage.
\item \textbf{P3 --- Existing-rule upgrades:} NOT IN та GROUP BY first; then ordering,
      SARGability, joins, DISTINCT.
\item \textbf{P4 --- New rules:} fan-out/chasm, join-key mismatch, scalar subquery.
\item \textbf{P5 --- Outputs/tests:} tri-state evidence, risk vector, DuckDB/reporting.
\item \textbf{P6 --- Pilot:} Spider Train development + held-out pilot.
\item \textbf{P7 --- Full study:} C0--C3, baselines, annotation.
\item \textbf{P8 --- Analysis/article/artifact:} statistics, frozen outputs, paper.
\end{enumerate}

%==================================================================
\section{План майбутньої статті}
%==================================================================

\subsection{Рекомендована назва}

\begin{thesisbox}
\begin{english}
\textbf{When Does Schema Context Help? A Controlled Study of Catalog-Grounded SQL
Auditing for Text-to-SQL Benchmarks}
\end{english}
\end{thesisbox}

Альтернативи:

\begin{itemize}
\item \begin{english}Beyond the AST: Catalog-Grounded Detection of Structural SQL Risks in
Text-to-SQL Benchmarks\end{english}
\item \begin{english}From Parse Trees to Catalog Evidence: Measuring the Value of
Schema Context in Gold-SQL Auditing\end{english}
\item \begin{english}Schema-Grounded Static Analysis of Text-to-SQL Reference Queries:
Precision, Abstention, and Cross-Dataset Evidence\end{english}
\end{itemize}

\subsection{Однореченнєва теза}

\begin{quote}
Catalog grounding materially changes which structural SQL findings are justified,
but its value depends on metadata completeness, provenance, and query lineage; a
reliable auditor must therefore measure both improved precision and principled
abstention.
\end{quote}

\subsection{Заплановані внески}

\begin{enumerate}
\item Context model із provenance/reliability і dual native/DDL acquisition.
\item Lineage-aware tri-state schema-conditioned rule engine.
\item Six refined UJIT families та три нові Text-to-SQL-relevant rules.
\item Controlled C0--C3 study на Spider/BIRD із Gretel DDL stress test.
\item Blind disagreement labels і fair comparison із \texttt{sqlsure}.
\item Reproducible artifact із evidence, manifests, rules і analysis scripts.
\end{enumerate}

\subsection{Abstract blueprint}

\begin{enumerate}
\item \textbf{Problem:} executable gold SQL може містити risks, яких AST-only або
      execution-only checks не розрізняють.
\item \textbf{Gap:} schema-aware systems існують, але incremental value, metadata
      reliability і abstention систематично не виміряні для reference-query audits.
\item \textbf{Method:} lineage-aware context, tri-state rules, risk vector, C0--C3.
\item \textbf{Study:} Spider/BIRD primary, Gretel DDL external validity, sqlsure baseline,
      blind human labels.
\item \textbf{Results:} вставити лише після freeze та evaluation: suppression,
      additions, unknown rates, precision, coverage, runtime.
\item \textbf{Implication:} benchmark QA має поєднувати AST, trustworthy catalog facts,
      execution і semantic review, не змішуючи їхні claims.
\end{enumerate}

\subsection{Section outline}

\begin{enumerate}
\item Introduction and contributions.
\item Background and relation to prior work.
\item Catalog-grounded analysis model.
\item Schema acquisition, provenance and lineage.
\item Rule specifications and risk model.
\item Study design.
\item Metadata reliability results.
\item C0--C3 ablation and human accuracy.
\item Comparison with sqlsure and other baselines.
\item Cross-dataset results and Gretel stress test.
\item Discussion, threats and implications.
\item Conclusion and artifact.
\end{enumerate}

\subsection{Заплановані figures і tables}

\begin{itemize}
\item Figure 1: AST + native/DDL catalog architecture.
\item Figure 2: scope/lineage propagation example.
\item Figure 3: decision flow \presentstate/\absentstate/\unknown.
\item Table 1: explicit delta from UJIT/e-Informatica/calibration.
\item Table 2: rule specifications та required facts.
\item Table 3: metadata coverage/reliability by corpus.
\item Table 4: C0--C3 migration matrix.
\item Table 5: blind precision/recall by rule family.
\item Table 6: matched overlap with sqlsure.
\item Table 7: runtime, cache and materialization cost.
\item Figure 4: schema-defect to affected-query graph.
\end{itemize}

%==================================================================
\section{Claims: що можна і чого не можна заявляти}
%==================================================================

\subsection{Допустимі claims після успішного experiment}

\begin{itemize}
\item schema context statistically significantly improves precision for specified
      rule families;
\item explicit abstention prevents unsupported conclusions under incomplete metadata;
\item provenance sources lead to measurably different rule decisions;
\item lineage increases applicable coverage on nested SQL;
\item proposed system complements sqlsure on non-overlapping nullability, FD, ordering,
      type та scalar-subquery risks;
\item DDL materialization enables scalable schema-conditioned auditing without
      populated databases, with clearly stated limits.
\end{itemize}

\subsection{Claims, яких слід уникати}

\begin{riskbox}
\begin{itemize}
\item «First schema-aware SQL antipattern detector».
\item «First schema-aware audit of Spider/BIRD gold SQL».
\item «Every flag is a wrong gold query».
\item «Missing FK proves a wrong join».
\item «Observed uniqueness proves a key».
\item «Execution success validates semantics».
\item «Gretel має 100\,000 повноцінних populated databases».
\item «sqlsure has zero false positives» без відтворення його exact protocol.
\item «Наш method dominates sqlsure» за глобальним score для різних taxonomies.
\item «Schema awareness solves NL--SQL intent mismatch».
\end{itemize}
\end{riskbox}

%==================================================================
\section{Threats to validity}
%==================================================================

\begin{longtable}{>{\raggedright\arraybackslash}p{4cm}
                  >{\raggedright\arraybackslash}p{9.8cm}}
\toprule
\textbf{Threat} & \textbf{Mitigation} \\
\midrule
\endfirsthead
\toprule
\textbf{Threat} & \textbf{Mitigation} \\
\midrule
\endhead
Missing constraints &
Absence $\rightarrow$ \unknown; provenance і coverage reporting. \\
Incorrect declarations &
Compare metadata/native catalog; consistency checks; schema-defect label. \\
SQLite permissive typing &
Coarse type families; portability axis; no universal correctness claim. \\
Current data cannot prove constraints &
Separate declared, data-consistent і observed-only facts. \\
Intent unknown &
Label structural proposition, not whole NL--SQL correctness. \\
Correlated queries per DB &
Clustered inference by \texttt{db\_id}. \\
Rule tuning on known cases &
Synthetic/Spider Train development; freeze; held-out BIRD/Spider evaluation. \\
Gretel duplication and heterogeneity &
DDL and canonical SQL hashes; separate reporting. \\
Baseline taxonomy mismatch &
Matched rules only; pin versions and parser coverage. \\
Reviewer bias &
Blind dual review and adjudication. \\
Publication overlap &
Explicit relation-to-prior-work section; no recycled headline results. \\
Rapidly changing related work &
Repeat literature and sqlsure status search immediately before submission. \\
\bottomrule
\caption{Основні threats і mitigations.}
\end{longtable}

%==================================================================
\section{Decision gates}
%==================================================================

\begin{description}[leftmargin=1.7cm,labelwidth=1.25cm]
\item[G1] Metadata coverage достатнє для заявлених rules. Інакше primary scope
звужується до Spider/BIRD.
\item[G2] Pilot NOT IN + join/cardinality показує precision/abstention improvement,
а не лише інші counts.
\item[G3] Resolver проходить nested/alias/composite-key fixture suite без silent guesses.
\item[G4] Є достатньо позитивів для human evaluation кожної headline family; рідкісні
classes подаються descriptive.
\item[G5] sqlsure comparison відтворюваний на pinned commit і однакових inputs.
\item[G6] Abstract не містить старих UJIT/e-Informatica чисел як нових results.
\item[G7] Перед submission перевірено, чи не з'явилась sqlsure paper/preprint.
\end{description}

\begin{warningbox}
Якщо G2 не проходить, це не автоматично провал. Можлива чесна negative-result study
або tool paper про те, що incomplete schema metadata дає менший виграш, ніж очікувалося.
Але не слід наперед писати сильний abstract без pilot evidence.
\end{warningbox}

%==================================================================
\section{Artifact і reproducibility checklist}
%==================================================================

\begin{multicols}{2}
\begin{itemize}
\item pinned code commit;
\item dataset/version manifest;
\item checksums;
\item DDL normalization version;
\item schema-context serialization;
\item rule specifications;
\item matched fixtures;
\item baseline wrappers;
\item sqlsure pinned snapshot;
\item annotation guideline;
\item blinded labels;
\item adjudication log;
\item statistical scripts;
\item raw migration matrices;
\item runtime environment;
\item generated reports;
\item known limitations;
\item relation-to-prior-work statement.
\end{itemize}
\end{multicols}

%==================================================================
\section{Практичний порядок наступних дій}
%==================================================================

\begin{enumerate}
\item Визначити canonical code repository/branch і freeze C0 behavior.
\item Зробити schema feasibility inventory для Spider, BIRD, Gretel.
\item Реалізувати context builder з native та DDL-materialized paths.
\item Реалізувати resolver до додавання rules.
\item Зробити matched pilot для NOT IN і join/cardinality.
\item Прогнати decision gates G1--G3.
\item Додати решту primary rules і outputs.
\item Freeze specifications і evaluation commit.
\item Запустити C0--C3 та pinned baselines.
\item Сформувати blind disagreement sample.
\item Провести review/adjudication і statistics.
\item Написати Results лише після завершення аналізу.
\item Випустити artifact разом із paper/preprint.
\end{enumerate}

%==================================================================
\section{Ключові джерела та посилання}
%==================================================================

\begin{enumerate}
\item Sabadosh, V. M., Kotsovsky, V. M. (2026).
  \href{https://doi.org/10.23939/ujit2026.01.135}
  {\textit{Automated Detection and Comparative Analysis of SQL Antipatterns in
  Text-to-SQL Datasets}}. Ukrainian Journal of Information Technology, 8(1), 135--143.
  \href{https://science.lpnu.ua/sites/default/files/journal-paper/2026/apr/42369/15sabadoshkocovskiy135-143.pdf}{PDF}.
\item Dintyala, P., Narechania, A., Arulraj, J. (2020).
  \href{https://doi.org/10.1145/3318464.3389754}
  {\textit{SQLCheck: Automated Detection and Diagnosis of SQL Anti-Patterns}}.
  ACM SIGMOD.
\item Arora, T. (2026).
  \href{https://github.com/sqlsure/sqlsure}{\texttt{sqlsure}: Semantic inspector for SQL},
  software artifact v0.1.1.
  \href{https://github.com/sqlsure/sqlsure/tree/d1b98e741f848ce7d756ee1d672e206d6ec39daa}
  {Frozen snapshot used in this plan}.
\item Yu, T., et al. (2018).
  \href{https://doi.org/10.18653/v1/D18-1425}
  {\textit{Spider: A Large-Scale Human-Labeled Dataset for Complex and Cross-Domain
  Semantic Parsing and Text-to-SQL Task}}.
\item Li, J., et al. (2023).
  \href{https://proceedings.neurips.cc/paper_files/paper/2023/hash/83fc8fab1710363050bbd1d4b8cc0021-Abstract-Datasets_and_Benchmarks.html}
  {\textit{Can LLM Already Serve as a Database Interface? A Big Bench for
  Large-Scale Database Grounded Text-to-SQLs}}.
\item Liu, X., et al. (2025).
  \href{https://doi.org/10.1145/3711896.3737427}
  {\textit{NL2SQL-BUGs: A Benchmark for Detecting Semantic Errors in NL2SQL
  Translation}}.
\item SQLGlot contributors.
  \href{https://github.com/tobymao/sqlglot}{\textit{SQLGlot: Python SQL Parser and
  Transpiler}}.
\item Gretel.ai.
  \href{https://huggingface.co/datasets/gretelai/synthetic_text_to_sql}
  {\textit{synthetic\_text\_to\_sql dataset}}.
\end{enumerate}

\vfill
\begin{center}
\footnotesize\sffamily\color{accent}
Окремий planning artifact. Не замінює наявний FullArticlePlan для
calibration/reconciliation study.
\end{center}

\end{document}
